\documentclass[../../../main]{subfiles}

\ifSubfilesClassLoaded{%
    \zexternaldocument*{../../1.Category-Theory/exported/aux/main}%
    \zexternaldocument*{../1.Additive-and-abelian-categories/exported/aux/main}%
}{}

\begin{document}
\section{Complexes and cohomology}

Throughout this chapter we will be working over a fixed \hyperref[def:abelian-category]{abelian category} \(\category{A}\).

\begin{definition}[Complex] \label{def:complex}
  A sequence of objects and morphisms
  \begin{center}
    \begin{tikzcd}
       \cdots \arrow[r, "d_{i+2}"'] & C_{i+1} \arrow[r, "d_{i+1}"'] & C_{i} \arrow[r, "d_{i}"'] & C_{i-1} \arrow[r, "d_{i-1}"'] & \cdots
    \end{tikzcd}
  \end{center}
  in \(\category{A}\) is a \textbf{chain complex} if for every \(i \in \mathbb{Z} \quad d_{i} \circ d_{i+1} = 0\).
  We denote a chain complex as \((C_{\bullet}, d_{\bullet})\) or simply as \(C_{\bullet}\).

  By reversing the order, we say that a sequence
  \begin{center}
    \begin{tikzcd}
      \cdots \arrow[r, "d^{i-2}"] & C^{i-1} \arrow[r, "d^{i-1}"] & C^{i} \arrow[r, "d^{i}"] & C^{i+1} \arrow[r, "d^{i+1}"] & \cdots
    \end{tikzcd}
  \end{center}
  is a \textbf{cochain complex} if for every \(i \in \mathbb{Z} \quad d^{i} \circ d^{i-1} = 0\).
  We denote a cochain complex as \((C^{\bullet}, d^{\bullet})\) or simply as \(C^{\bullet}\).
\end{definition}

\begin{remark}
  In purelly algebraic terms, chains are cochains are \textit{equivalent}.
  We will adopt the cochain notation, but all proofs and definitions translate to the chain notation.
  Therefore from now onwards we will be refering to \textit{cochains} as \textit{chains} or \textit{sequences}.
\end{remark}

\begin{definition}[Morphism of complexes] \label{def:complex-morphism}
  Let \(C^{\bullet}, D^{\bullet}\) be two \hyperref[def:complex]{complexes} in \(\category{A}\).
  A \textbf{morphism of complexes} \(f: C^{\bullet} \to D^{\bullet}\) is a collection of morphisms \(\{f^{i}: C^{i} \to D^{i}\}_{i \in \mathbb{Z}}\) making, for every \(i \in \mathbb{Z}\), the following diagram commute
  \begin{center}
    \begin{tikzcd}
      C^{i} \arrow[r, "d_{C}^{i}"] \arrow[d, "f^{i}"'] & C^{i+1} \arrow[d, "f^{i+1}"] \\
      D^{i} \arrow[r, "d_{D}^{i}"] & D^{i+1}
    \end{tikzcd}
  \end{center}
\end{definition}

\begin{definition} \label{def:complex-category}
  Given any \hyperref[def:abelian-category]{abelian category} \(\category{A}\) we can construct a \textbf{complex category} \(\Complex{\category{A}}\) by taking
  \(\operatorname{Ob}_{\Complex{\category{A}}}\) to be the collection of all \hyperref[def:complex]{complexes} with values in \(\category{A}\) and, for every \(C^{\bullet}, D^{\bullet} \in \operatorname{Ob}_{\Complex{\category{A}}}\), \(\Hom_{\Complex{\category{A}}}(C^{\bullet}, D^{\bullet})\) to be the collection of all \hyperref[def:complex-morphism]{morphisms of complexes} between \(C^{\bullet}\) and \(D^{\bullet}\).
\end{definition}

\begin{remark} \label{remark:mono-epi-in-complexes}
  A \hyperref[def:complex-morphism]{morphism of complexes} \(f: C^{\bullet} \to D^{\bullet}\) is a \hyperref[def:mono-epi-morphism]{monomorphism} (resp. \hyperref[def:mono-epi-morphism]{epimorphism}) if for every \(i \in \mathbb{Z}\) the morphism \(f^{i}: C^{i} \to D^{i}\) is a \hyperref[def:mono-epi-morphism]{monomorphism} (resp. \hyperref[def:mono-epi-morphism]{epimorphism}) in \(\category{A}\).
\end{remark}

\begin{theorem} \label{thm:complex-category-is-abelian}
  If \(\category{A}\) is an \hyperref[def:abelian-category]{abelian-category}, so is its \hyperref[def:complex-category]{category of complexes} \(\Complex{\category{A}}\).
\end{theorem}
\begin{proof}
  By taking addition \textit{elementwise} and \(0^{\bullet}\) as the \hyperref[def:zero-object]{zero object} we can make \(\Complex{\category{A}}\) into an \hyperref[def:additive-category]{additive category}.
  Therefore we only need to prove that \(\Complex{\category{A}}\) verifies the axioms of an \hyperref[def:abelian-category]{abelian category}.

  Let \(f: A^{\bullet} \to B^{\bullet}\) be a \hyperref[def:complex-morphism]{morphism of complexes}.

  Since \(\category{A}\) is \hyperref[def:abelian-category]{abelian}, for every \(i \in \mathbb{Z}\) the map \(f^{i}: A^{i} \to B^{i}\) admits a \hyperref[def:categorical-kernel]{kernel} \(\ker f^{i}: K^{i} \monoarrow A^{i}\) and a \hyperref[def:categorical-cokernel]{cokernel} \(\coker f^{i}: B^{i} \epiarrow C^{i}\).

  Consider the diagram
  \begin{center}
    \begin{tikzcd}
      K^{i} \arrow[d, tail, "\ker f^{i}"'] & K^{i+1} \arrow[d, tail, "\ker f^{i+1}"] \\
      A^{i} \arrow[r, "d_{C}^{i}"] \arrow[d, "f^{i}"'] & A^{i+1} \arrow[d, "f^{i+1}"] \\
      B^{i} \arrow[r, "d_{D}^{i}"] \arrow[d, twoheadrightarrow, "\coker f^{i}"'] & B^{i+1} \arrow[d, twoheadrightarrow, "\coker f^{i+1}"] \\
      C^{i}  & C^{i+1}  \\
    \end{tikzcd}
  \end{center}
  Since the inner square is commutative we get unique morphisms \(K^{i} \to K^{i+1}\) and \(C^{i} \to C^{i+1}\) making the whole diagram commute, obtaining \hyperref[def:complexes]{chain complexes} \(K^{\bullet}\) and \(C^{\bullet}\) and \hyperref[def:complex-morphism]{morphisms of chains} \(\ker f^{\bullet}, \coker f^{\bullet}\).
  It is then easy to check that \(\ker f^{\bullet}\) and \(\coker f^{\bullet}\) verify the corresponding universal properties.

  On the other hand, by \autoref{remark:mono-epi-in-complexes} the other axiom is easily verified \textit{element by element}.
\end{proof}


\begin{definition}[Bounded complexes] \label{def:bounded-complexes}
  Let \(C^{\bullet}\) be a \hyperref[def:complex]{complex}. We say that \(C^{\bullet}\) is
  \begin{itemize}
    \item \textbf{bounded below} if there exists \(k \in \mathbb{Z}\) such that \(\forall i \leq k \quad C^{i} = 0\).
    \item \textbf{bounded above} if there exists \(k \in \mathbb{Z}\) such that \(\forall i \geq k \quad C^{i} = 0\).
    \item \textbf{bounded} if it is both bounded below and above.
  \end{itemize}
\end{definition}

\begin{notation}
  We denote by \(\Complex[-, +, b]{\category{A}} \hookrightarrow \Complex{\category{A}}\) the \hyperref[def:full-subcategory]{full subcategory} of the \hyperref[def:complex-category]{chain complex category} consisting of \hyperref[def:bounded-complexes]{bounded below}, \hyperref[def:bounded-complexes]{bounded above} and \hyperref[def:bounded-complexes]{bounded} complexes respectively.

  We also have variants \(\Complex[\geq 0, \leq 0]{\category{A}} \hookrightarrow \Complex{\category{A}}\) of complexes with \(\forall i \geq 0\) (resp \(\forall  i \leq 0\)) \(C^{i} = 0\).

  By \autoref{thm:complex-category-is-abelian} all of this categories are \hyperref[def:abelian-category]{abelian}.
\end{notation}


\begin{remark} \label{remark:inclussion-into-complex}
  For all integers \(i \in \mathbb{Z}\) there exists a \hyperref[def:functor-types]{fully faithfull functor}
  \begin{align*}
    \iota_{r}: \category{A} \to \Complex{\category{A}}
  \end{align*}
  sending an object \(A \in \category{A}\) to a \hyperref[def:complex]{complex}
  \begin{center}
    \begin{tikzcd}
      \cdots \arrow[r] & 0 \arrow[r, "d^{r-1}"] & A \arrow[r, "d^{r}"] & 0 \arrow[r] & \cdots
    \end{tikzcd}
  \end{center}
  and a morphism \(\varphi: A \to B\) in \(\category{A}\) to a \hyperref[def:complex-morphism]{morphism of complexes}
  \begin{center}
    \begin{tikzcd}
      \cdots \arrow[r] & 0 \arrow[r, "d^{r-1}"] \arrow[d] & A \arrow[r, "d^{r}"] \arrow[d, "\varphi"] & 0 \arrow[r] \arrow[d] & \cdots \\
      \cdots \arrow[r] & 0 \arrow[r, "d^{r-1}"] & B \arrow[r, "d^{r}"] & 0 \arrow[r] & \cdots
    \end{tikzcd}
  \end{center}

  In this way, we can identify \(\category{A}\) as its \textit{image} via \(\iota = \iota_{0}\) in \(\Complex{\category{A}}\).
\end{remark}

\begin{definition}[Shift functor] \label{def:shift-functor}
  For any integer \(i \in \mathbb{Z}\) we define the \textbf{shift} \hyperref[def:functor]{functor}
  \begin{align*}
    \Complex{\category{A}} &\to \Complex{\category{A}} \\
    M^{\bullet} &\mapsto M[i]^{\bullet}
  \end{align*}
  by \(M[i]^{k} = M^{k+i}\) and \(d_{M[i]^{\bullet}}^{k} = (-1)^{i}d_{M^{\bullet}}^{k + i}\).
\end{definition}

\begin{definition}[Exact complex] \label{def:exact-complex}
  A \hyperref[def:complex]{complex} \(C^{\bullet}\) is \textbf{exact} at position \(i\) if the sequence
  \begin{center}
    \begin{tikzcd}
      C^{i-1} \arrow[r, "d^{i-1}"] & C^{i} \arrow[r, "d^{i}"] & C^{i+1}
    \end{tikzcd}
  \end{center}
  is \hyperref[def:exact-sequence]{exact}.

  If the \hyperref[def:complex]{complex} is exact at every position, we say that it is \textbf{exact}.
\end{definition}



\end{document}
